\documentclass[a4paper]{article}
% Español (México) con XeLaTeX: fontspec para fuentes Unicode, polyglossia
% para la división silábica y los nombres del documento en la variante mexicana.
\usepackage{fontspec}
\usepackage{polyglossia}
\setdefaultlanguage[variant=mexican]{spanish}
% Los nombres chinos van como «pinyin (original)»: xeCJK da a los caracteres
% Han su propia fuente; sin ella XeLaTeX los omite con un aviso.
% FandolSong no cubre algunos caracteres de nombres (U+73FA); WenQuanYi Zen Hei sí.
\usepackage[AutoFallBack=true]{xeCJK}
\setCJKmainfont{FandolSong-Regular.otf}
\setCJKfallbackfamilyfont{\CJKrmdefault}{WenQuanYi Zen Hei}
% polyglossia no traduce los nombres de listings.
\AtBeginDocument{\providecommand{\lstlistingname}{}\renewcommand{\lstlistingname}{Listado}%
  \providecommand{\lstlistlistingname}{}\renewcommand{\lstlistlistingname}{Índice de listados}}
\usepackage{amsmath, amssymb}
\usepackage{graphicx}
\usepackage[margin=2.5cm]{geometry}
\usepackage[most]{tcolorbox}
\usepackage{etoolbox}
\usepackage{listings}
\usepackage{hyperref}
\usepackage{booktabs}
\usepackage{subcaption}
\usepackage{float}

\newtcolorbox{knowledgebox}[1]{enhanced, colback=blue!5!white, colframe=blue!75!black, colbacktitle=blue!75!black, coltitle=white, fonttitle=\bfseries, title={#1}, attach boxed title to top left={yshift=-2mm, xshift=2mm}, boxrule=1pt, sharp corners}
\newtcolorbox{importantbox}[1]{enhanced, colback=yellow!10!white, colframe=yellow!80!black, colbacktitle=yellow!80!black, coltitle=black, fonttitle=\bfseries, title={#1}, sharp corners}
\newtcolorbox{warningbox}[1]{enhanced, colback=red!5!white, colframe=red!75!black, colbacktitle=red!75!black, coltitle=white, fonttitle=\bfseries, title={#1}, sharp corners}

\lstset{language=Python, basicstyle=\ttfamily\small, keywordstyle=\color{blue}, stringstyle=\color{red!60!black}, commentstyle=\color{green!60!black}, breaklines=true, frame=single, numbers=left, numberstyle=\tiny\color{gray}, captionpos=b, extendedchars=true}

\newcommand{\notetitle}{[LLM Agents SP25] Multimodal Agents: Perception to Action}
\newcommand{\noteauthors}{Elaboradas a partir de la clase de Caiming Xiong}
\newcommand{\notedate}{2026-04-02}
\newcommand{\videochannel}{Berkeley RDI}
\newcommand{\videopublishdate}{2025-03-20}
\newcommand{\videoduration}{1:28:21}
\newcommand{\videourl}{https://www.youtube.com/watch?v=n\_\_Tim8K2IY}
\newcommand{\videocoverpath}{cover.jpg}
\newcommand{\repourl}{https://github.com/hqhq1025/ai-course-notes}

\usepackage{fancyhdr}
\pagestyle{fancy}
\fancyhf{}
\fancyfoot[C]{\thepage}
\fancyfoot[R]{\footnotesize\color{gray}\href{\repourl}{AI Course Notes}}
\renewcommand{\headrulewidth}{0pt}
\renewcommand{\footrulewidth}{0pt}

\begin{document}
\begin{titlepage}
\centering
{\Large Notas del curso\par}
\vspace{1.2cm}
{\huge\bfseries \notetitle\par}
\vspace{0.8cm}
{\large \noteauthors\par}
\vspace{0.3cm}
{\large \notedate\par}
\vspace{1.2cm}
\ifdefempty{\videocoverpath}{}{\includegraphics[width=0.82\textwidth,height=0.45\textheight,keepaspectratio]{\videocoverpath}\par}
\vfill
\begin{tcolorbox}[width=0.9\textwidth, colback=black!2!white, colframe=black!60, sharp corners]
\textbf{Autor o canal del video}: \videochannel\par
\textbf{Fecha de publicación}: \videopublishdate\par
\textbf{Duración del video}: \videoduration\par
\ifdefempty{\videourl}{}{\textbf{Enlace del video}: \href{\videourl}{\nolinkurl{\videourl}}\par}
\textbf{Página del proyecto}: \href{\repourl}{\nolinkurl{\repourl}}\par
\end{tcolorbox}
\end{titlepage}

\tableofcontents
\newpage

\section{De la saturación de los puntos de referencia a los Agentes multimodales: por qué se reescribió el espacio del problema}

\subsection{Hilo conductor de la charla y juicio central}
Esta charla la dio Caiming Xiong, de Salesforce AI Research, y su hilo conductor es muy claro: los puntos de referencia tradicionales de tareas de texto se acercaron a la saturación en los últimos años, y el desafío de la siguiente etapa ya no es «responder preguntas fuera de línea», sino «completar tareas de manera continua en un entorno digital real». Esta definición desplaza el objeto de estudio de tareas de PLN estáticas hacia sistemas de Agentes interactivos, de varios pasos y que operan a través de aplicaciones.

\begin{importantbox}{De «saber responder» a «saber hacer»}
En el paradigma de los Agentes multimodales, el valor de un modelo ya no se mide solo por si una frase que produce es correcta, sino por si puede completar de manera estable un flujo de trabajo real: leer la interfaz, entender el estado, planear la acción, ejecutar la operación, verificar el resultado y volver a iterar. Este ciclo cerrado es el verdadero umbral en un entorno de producción.
\end{importantbox}

\subsection{Por qué «multimodal» ya no es opcional}
El expositor define al Agente multimodal como un sistema Vision-Language-Action (VLA). La razón es directa: el flujo de trabajo real no es una API de puro texto, sino sobre todo interfaces gráficas, páginas web, software de escritorio, interfaces móviles e incluso interacción con periféricos. Solo si el Agente cuenta a la vez con comprensión visual, razonamiento en lenguaje y generación de acciones puede llegar a hacerse cargo de un proceso complejo.

\begin{knowledgebox}{Comparación de entradas y salidas del Agente multimodal}
\begin{itemize}
    \item La entrada pasa de ser «texto de un solo turno» a ser «instrucción del usuario + captura de pantalla + árbol de accesibilidad + estado del entorno»;
    \item la salida pasa de ser «respuesta en lenguaje natural» a ser «secuencia de acciones ejecutables»;
    \item la evaluación pasa de ser «exactitud» a ser «tasa de finalización de la tarea, calidad de la trayectoria, eficiencia de los pasos y robustez».
\end{itemize}
\end{knowledgebox}

\begin{figure}[H]
\centering
\includegraphics[width=0.88\textwidth]{frame-01-benchmark.jpg}
\caption{Video, aproximadamente 00:01:30: el expositor muestra que los puntos de referencia recientes se saturan con rapidez, y a partir de ahí plantea que la evaluación de Agentes debe pasar de un banco de preguntas fuera de línea a un entorno real}
\end{figure}

\begin{warningbox}{Malentendido frecuente: confundir «saber usar herramientas» con «Agente listo para producción»}
Un éxito en una demostración no significa que el sistema esté listo para salir a producción. Un Agente de nivel de producción se preocupa por la estabilidad a largo plazo: recuperación de errores, manejo de interrupciones de tareas, consistencia del estado entre aplicaciones, límites de permisos y capacidad de auditoría. Ninguno de estos aspectos puede cubrirse con la puntuación de un solo benchmark.
\end{warningbox}

\subsection{Resumen de la sección}
El foco de la investigación sobre Agentes multimodales pasó de la «capacidad de responder» del modelo a la «capacidad de interactuar con el entorno». Este giro determina que, de ahora en adelante, deban construirse a la vez tres capas de base: un benchmark de entorno real, un pipeline de datos escalable y una arquitectura de modelo ejecutable.
\section{Mapa de aplicaciones: cómo los flujos de trabajo digitales se convierten en el principal campo de batalla del agente}

\subsection{Forma de la tarea: de la operación en una sola aplicación a la orquestación entre aplicaciones}
El docente enfatiza que las tareas reales de los usuarios normalmente no son una página con un solo clic, sino un proceso que atraviesa varias herramientas. Por ejemplo: buscar información en el navegador, organizar datos en una hoja de cálculo, enviar los resultados por correo electrónico, archivar el registro en un documento. Un agente multimodal necesita manejar un estado persistente y el contexto de la tarea, no solo responder a instrucciones de corto plazo.

\begin{importantbox}{Cuatro características de las tareas del mundo real}
\begin{enumerate}
    \item \textbf{Cadena larga}: a menudo supera los 20--50 pasos;
    \item \textbf{Entre interfaces}: navegador, escritorio, terminal y herramientas de Office aparecen mezclados;
    \item \textbf{No determinismo}: la carga de páginas, las ventanas emergentes y las fluctuaciones de red provocan cambios de estado;
    \item \textbf{Objetivo abierto}: muchas tareas admiten varias rutas viables, no una única respuesta estándar.
\end{enumerate}
\end{importantbox}

\subsection{Estratificación de las capacidades del agente: percepción, planeación, ejecución, verificación}
El curso descompone al agente en una pila de capacidades, en lugar de tratarlo como un modelo de caja negra. La capa de percepción se encarga de entender la interfaz y los objetos semánticos; la capa de planeación decide el siguiente subobjetivo; la capa de ejecución produce acciones ejecutables; la capa de verificación determina si la tarea se completó y corrige la trayectoria. Esta estratificación permite ubicar mejor los cuellos de botella del sistema.

\begin{knowledgebox}{Por qué el agente multimodal necesita ``visibilidad del proceso''}
Cuando una tarea falla, el equipo debe saber si el problema está en la percepción, la planeación o la ejecución. Si todos los fallos se atribuyen a que ``el modelo no es lo bastante fuerte'', no es posible formular una estrategia de optimización eficaz. La observabilidad es el núcleo de la eficiencia de iteración del agente.
\end{knowledgebox}

\subsection{Prioridad a los escenarios digitales, los escenarios físicos después}
Desde el punto de vista del ritmo de implementación, el docente coloca claramente los entornos digitales (web, escritorio, móvil) como la dirección de mayor prioridad. La razón es que el ciclo de retroalimentación es más rápido, los guiones de evaluación son más fáciles de construir y el costo es más controlable. Los agentes en el mundo físico son importantes, pero su infraestructura de datos y evaluación es más pesada y su velocidad de iteración es más lenta.

\begin{warningbox}{No confundir prematuramente la ruta de la robótica con la del agente de GUI}
Ambos son agentes, pero sus restricciones de ingeniería son completamente distintas. El cuello de botella del agente de GUI está en la comprensión de la interfaz y la orquestación de tareas; el agente físico añade además el error de percepción, las restricciones de seguridad y la estabilidad del actuador. Discutirlos juntos oculta la naturaleza del problema.
\end{warningbox}

\subsection{Resumen de la sección}
El curso define el problema de las aplicaciones como una tarea de largo plazo que es ``entre aplicaciones, de múltiples pasos y recuperable''. De ahí se desprende el requisito de ingeniería: primero cerrar a fondo el ciclo de las tareas digitales, y después expandirse a escenarios más complejos.

\section{Actualización del sistema de benchmarks: de los benchmarks web a OS-World}

\subsection{El valor y los límites de los benchmarks existentes}
El expositor repasó benchmarks como Mind2Web, WebArena y VisualWebArena. Estos trabajos impulsaron la investigación en Web Agent, pero su cobertura de entornos, tipos de tarea y escalabilidad siguen sin bastar para representar escenarios reales de ``computer use'', en particular las tareas entre aplicaciones y de objetivo abierto.

\begin{knowledgebox}{Limitaciones típicas}
\begin{itemize}
    \item La distribución de tareas se inclina hacia la navegación web, con una cobertura estrecha;
    \item Muchas tareas siguen rutas implícitamente plantilladas, fáciles de sobreajustar con una política;
    \item El soporte para aplicaciones de escritorio, sistemas de archivos y flujos entre programas es insuficiente;
    \item Es difícil sostener un ciclo integrado de ``desarrollo, prueba y evaluación''.
\end{itemize}
\end{knowledgebox}

\subsection{Los objetivos de diseño de OS-World}
OS-World se posiciona como un entorno informático real más general: dentro de una máquina virtual ofrece un escritorio interactivo y un ecosistema de aplicaciones, y permite que los investigadores definan tareas, ejecuten agentes, recolecten trayectorias y corran scripts de evaluación, dejando además espacio para ampliarlo.

\begin{importantbox}{Tres propiedades clave de OS-World}
\begin{enumerate}
    \item \textbf{Realismo}: se acerca al comportamiento real de un sistema operativo y sus aplicaciones;
    \item \textbf{Escalabilidad}: permite agregar de manera sostenida nuevas tareas y combinaciones de aplicaciones;
    \item \textbf{Evaluabilidad}: cada tarea puede vincularse a una lógica de verificación programada.
\end{enumerate}
\end{importantbox}

\begin{figure}[H]
\centering
\includegraphics[width=0.88\textwidth]{frame-02-osworld.jpg}
\caption{Video aproximadamente en 00:10:32: se propone OS-World como un entorno real y escalable para el desarrollo y la evaluación de agentes multimodales}
\end{figure}

\subsection{Perspectiva comparativa: por qué OS-World se acerca más a los problemas de producción}
\begin{table}[H]
\centering
\begin{tabular}{p{3.2cm}p{3.7cm}p{3.8cm}}
\toprule
\textbf{Dimensión} & \textbf{Benchmarks web tradicionales} & \textbf{OS-World} \\
\midrule
Cobertura de entornos & Principalmente el dominio web & Web + aplicaciones de escritorio + interacción con el sistema \\
Forma de las tareas & Se inclina hacia navegación y formularios & Flujos entre aplicaciones, de varios pasos y de objetivo abierto \\
Método de evaluación & Orientado al resultado de la página & Verificación por script + recompensa por cumplimiento parcial \\
Capacidad de ampliación & Costo alto para construir tareas nuevas & Soporta agregar tareas y configuraciones de manera sistemática \\
Adaptación de ingeniería & Se inclina hacia la evaluación fuera de línea & Desarrollo, entrenamiento y evaluación integrados \\
\bottomrule
\end{tabular}
\caption{OS-World no busca sustituir todos los benchmarks, sino llenar el vacío central de las ``tareas informáticas reales''}
\end{table}

\begin{warningbox}{Riesgo de traslado de la evaluación}
Si un equipo solo optimiza sobre un único benchmark, el modelo tiende a desarrollar una ``especulación de política'' frente al script de evaluación. El curso recomienda hacer validación cruzada en múltiples entornos y familias de tareas para evitar el sobreajuste al benchmark.
\end{warningbox}

\subsection{Resumen de la sección}
El significado de OS-World no es simplemente sumar un benchmark más, sino desplazar la investigación de agentes de la ``competencia de bancos de preguntas'' hacia la ``ingeniería de entornos''. Esto ofrece una base experimental unificada para el diseño de datos y modelos posterior.
\section{Mecanismo de ejecución de OS-World: observación, acción y ciclo de interacción}

\subsection{Diseño del espacio de observación}
El curso explicó entradas de observación de múltiples fuentes: capturas de pantalla, el árbol de accesibilidad (accessibility tree), la salida de terminal y flujos de estado personalizados. Cada Agent puede elegir, según sus necesidades, una combinación de entradas dominada por lo visual o por lo textual.

\begin{knowledgebox}{Por qué conviene conservar la captura de pantalla y el árbol de accesibilidad a la vez}
La captura de pantalla ofrece un contexto visual integral, pero con poca estructura; el árbol de accesibilidad cuenta con jerarquía de objetos e información de atributos, lo que facilita la localización precisa y las operaciones interpretables. La complementariedad de ambos mejora de manera notable la estabilidad en la generación de acciones.
\end{knowledgebox}

\begin{figure}[H]
\centering
\includegraphics[width=0.88\textwidth]{frame-03-observation.jpg}
\caption{Video aproximadamente en 00:16:30: el ponente muestra la forma de observación combinada de screenshot y accessibility tree}
\end{figure}

\subsection{Espacio de acción e interfaz de ejecución}
Las acciones se expresan mediante instrucciones ejecutables que abarcan clics, entrada de texto, desplazamiento, atajos de teclado y operaciones combinadas. El Agent no genera directamente una ``sugerencia'', sino código de comportamiento ejecutable en la máquina virtual.

\begin{importantbox}{El valor de estandarizar el formato de las acciones}
Un formato de acción unificado sirve a la vez a tres fines: la alineación de los datos de entrenamiento, la operación estable del ejecutor y la reproducibilidad de los scripts de evaluación. Sin una capa de acciones estándar, tanto la comparación entre modelos como la acumulación de datos a gran escala fracasarían.
\end{importantbox}

\subsection{Ciclo de interacción}
El ciclo del sistema es ``recibir observación $\rightarrow$ planear acción $\rightarrow$ ejecutar en el entorno $\rightarrow$ nueva observación''. La tarea termina al alcanzar el objetivo, al activarse una condición de fallo o al llegar al número máximo de pasos. Este ciclo amplía el objetivo de entrenamiento del Agent, de la predicción de un solo paso a la toma de decisiones secuencial.

\begin{lstlisting}[caption={Ciclo de interacción simplificado de un Agent multimodal en OS-World}]
obs = env.reset(task_config)
for step in range(max_steps):
    action = agent.predict(obs, instruction)
    obs, reward, done, info = env.step(action)
    if done:
        break
\end{lstlisting}

\subsection{Detalles de ingeniería: seguridad y reproducibilidad}
El curso subrayó la importancia del aislamiento en máquinas virtuales y de los archivos de configuración de tareas. El entorno de investigación debe poder reiniciarse, registrarse y reproducirse, de modo que sea posible hacer experimentos comparativos estables y diagnosticar trayectorias de fallo.

\begin{warningbox}{Puntos de fallo comunes}
En un entorno real, lo que más se pasa por alto es la ``deriva de estado'': la caché, el estado de inicio de sesión, las ventanas emergentes y los procesos en segundo plano provocan que la misma tarea se comporte de manera inconsistente entre distintas ejecuciones. Es necesario reducir el efecto de esta deriva mediante instantáneas del entorno, scripts de inicialización y el registro de trazas.
\end{warningbox}

\subsection{Resumen de la sección}
Lo esencial de OS-World no es la captura de pantalla de la interfaz en sí, sino el diseño de ciclo cerrado de ``observación estructurada + acciones estándar + ejecución reproducible''. Esto lleva la investigación de Agents multimodales a una etapa de iteración ingenieril.
\section{Evaluación y recompensa: cómo convertir tareas abiertas en objetivos entrenables}

\subsection{Evaluación basada en ejecución (Execution-based Evaluation)}
El curso utiliza evaluación por ejecución de scripts, en lugar de calificación subjetiva humana. Una vez terminada la tarea, un script de verificación revisa archivos, el estado de la página, el contenido de la salida o variables del sistema, y así determina el grado de finalización.

\begin{importantbox}{Las tareas abiertas también admiten una evaluación computable}
Incluso si la tarea admite múltiples rutas, basta con que el estado del resultado sea verificable para poder construir una evaluación automática. Lo importante es traducir la «intención de la tarea» en «condiciones verificables», por ejemplo, la existencia de un archivo, el valor de un campo, el estado de un elemento de la página o el valor de retorno de un comando.
\end{importantbox}

\begin{figure}[H]
\centering
\includegraphics[width=0.88\textwidth]{frame-04-reward.jpg}
\caption{Video aproximadamente en 00:19:35: el expositor presenta la recompensa basada en ejecución y la lógica de calificación por finalización parcial}
\end{figure}

\subsection{Recompensa por finalización parcial y estratificación de dificultad}
El curso enfatiza que no debe usarse únicamente un resultado binario de éxito o fracaso. En tareas complejas, una puntuación de finalización parcial refleja el progreso con mayor granularidad y mejora de manera notable el problema de las señales de entrenamiento dispersas. En la práctica, esto debe combinarse con una estratificación de dificultad, para evitar que el modelo se quede estancado durante mucho tiempo en subtareas de nivel bajo.

\begin{knowledgebox}{Recomendaciones para el diseño de la función de recompensa}
\begin{itemize}
    \item Dividir el objetivo en hitos clave y establecer recompensas por segmentos;
    \item Aumentar la penalización de las acciones destructivas, para restringir las conductas riesgosas;
    \item Aplicar un tratamiento de decaimiento a las operaciones repetitivas sin efecto, para inhibir las trayectorias cíclicas;
    \item Gestionar el versionado de los scripts de evaluación, para garantizar la reproducibilidad.
\end{itemize}
\end{knowledgebox}

\subsection{Requisitos de robustez del sistema de evaluación}
El propio script de evaluación también puede ser «atacado» o eludido. El curso recomienda incorporar en la capa del script ramas para casos anómalos, verificaciones de límites y restricciones contra trampas, para evitar que el modelo aprenda estrategias oportunistas ajenas al objetivo de la tarea.

\begin{warningbox}{Reward Hacking es un riesgo real}
Si la recompensa solo cubre un estado local, el Agent puede aprender a «hacer trampa con la puntuación» en lugar de «completar la tarea». Por ejemplo, generando un falso estado de finalización mediante el caché de la interfaz, o aprovechando una vulnerabilidad del script para evadir pasos clave. El diseño de la evaluación debe preceder a la iteración del modelo.
\end{warningbox}

\subsection{Resumen de la sección}
El límite superior del entrenamiento del Agent está determinado en gran medida por el sistema de evaluación. La evaluación basada en ejecución, la recompensa parcial y la robustez de los scripts determinan en conjunto si el modelo puede progresar de manera estable en tareas abiertas.

\section{Cuello de botella de datos: de la anotación manual a la generación de trayectorias impulsada por tutoriales}

\subsection{Por qué los datos de trayectorias son el recurso escaso central}
El presentador señala que las trayectorias de interacción de varios pasos son mucho más costosas que los datos de instrucción de un solo turno. Cada trayectoria de alta calidad suele contener cambios de estado complejos, llamadas a herramientas y recuperación de errores; la recolección puramente manual difícil de sostener el entrenamiento a escala.

\begin{importantbox}{La escasez de datos no está en la ``cantidad de texto'', sino en la ``calidad de la trayectoria ejecutable''}
El valor de los datos de agentes proviene de tres puntos: contexto de tarea completo, acción ejecutable y resultado verificable. Si falta cualquiera de ellos, la ganancia de entrenamiento cae de manera considerable.
\end{importantbox}

\subsection{Construcción de datos impulsada por tutoriales (Tutorial-driven)}
El curso propone extraer plantillas de tareas de alto valor a partir de tutoriales de internet: primero se hace la recolección de tutoriales y el análisis estructurado, luego se convierten en descripciones de tareas estandarizadas y candidatos de pasos, y por último el agente los ejecuta en el entorno para generar las trayectorias.

\begin{figure}[H]
\centering
\includegraphics[width=0.88\textwidth]{frame-05-tutorial-pipeline.jpg}
\caption{Video aproximadamente en 00:44:40: pipeline de construcción de datos impulsada por tutoriales, que incluye extracción, estructuración, ejecución y filtrado}
\end{figure}

\begin{knowledgebox}{Ventajas del esquema impulsado por tutoriales}
\begin{itemize}
    \item Alta diversidad de tareas: cubre la distribución de problemas reales de los usuarios;
    \item Buena transferibilidad: tutoriales del mismo tipo pueden mapearse a tareas de distintos entornos;
    \item Costo controlable: reduce la anotación manual paso a paso;
    \item Fácil de escalar: puede integrarse con scripts de evaluación automática para formar un ciclo cerrado.
\end{itemize}
\end{knowledgebox}

\subsection{Control de calidad: de ``ejecutable'' a ``aprendible''}
El texto de un tutorial no equivale de manera natural a datos de entrenamiento. El curso enfatiza el filtrado en varios niveles: verificación de la ejecutabilidad de la tarea, verificación de la consistencia de la trayectoria, eliminación de pasos erróneos y reducción del peso de muestras con poca información. Los datos con mucho ruido bajan de manera considerable la eficiencia del entrenamiento.

\begin{warningbox}{No confundas ``recolección exitosa'' con ``dato utilizable''}
Muchas trayectorias, aunque se ejecutan por completo, casi no aportan ganancia al aprendizaje del modelo; por ejemplo, rutas de plantilla repetidas, ausencia de pasos de decisión clave o desalineación entre la acción y la observación. La gobernanza de datos debe integrarse al flujo principal de entrenamiento y no ser un parche fuera de línea.
\end{warningbox}

\subsection{Resumen de la sección}
La ruta de datos que propone el curso es ``recolección automática + gobernanza estructurada + verificación por script''. Esta ruta no busca la pureza absoluta, sino aumentar de manera sostenida la densidad de muestras efectivas con un costo controlado.
\section{Datos sintéticos y Action Calling: construcción de un volante de entrenamiento escalable}

\subsection{De las tareas de percepción a las tareas de acción}
El expositor enfatiza que la comprensión visual por sí sola (entender la interfaz) no basta para formar un Agente robusto; la capacidad debe transferirse a «invocar acciones». Por ello, el curso introduce un pipeline de datos sintéticos que convierte preguntas y respuestas sobre imágenes, OCR, localización y razonamiento de flujos en muestras ejecutables como acciones.

\begin{importantbox}{La esencia de Action Calling}
Action Calling no es una función adicional, sino el cambio del espacio de salida del modelo de «describir» a «ejecutar». Esto exige que los datos de entrenamiento incluyan a la vez intención, ruta de razonamiento, formato de la acción y retroalimentación de la ejecución.
\end{importantbox}

\begin{figure}[H]
\centering
\includegraphics[width=0.88\textwidth]{frame-06-synthetic-data.jpg}
\caption{Video aproximadamente en 01:02:32: el pipeline de datos sintéticos unifica preguntas visuales, cadenas de razonamiento e invocación de acciones en un mismo flujo de entrenamiento}
\end{figure}

\subsection{Flujo sintético típico}
\begin{enumerate}
    \item Muestrear escenarios de tarea y estados de la interfaz;
    \item Construir preguntas y objetivos (que incluyan localización, lectura, operación y verificación);
    \item Usar el modelo para generar candidatos de razonamiento y secuencias de acciones;
    \item Filtrar muestras de alta calidad con ayuda del ejecutor y del evaluador;
    \item Pasar a un entrenamiento por etapas y retroalimentar continuamente las muestras fallidas.
\end{enumerate}

\begin{knowledgebox}{Por qué funciona el entrenamiento por etapas}
En el curso se menciona varias veces: «primero aprender a ver, luego aprender a hacer, y después aprender la estrategia de largo plazo». Si desde el inicio se aprende de extremo a extremo una acción de cadena larga, el modelo cae con mayor facilidad en alta varianza y óptimos locales. Dividir el entrenamiento en etapas reduce la dificultad de la optimización y mejora el aprovechamiento de las muestras.
\end{knowledgebox}

\subsection{Los límites de los datos sintéticos}
Aunque las muestras sintéticas permiten escalar, el sesgo de distribución es inevitable. El curso recomienda mantener equilibrada la proporción entre trayectorias reales y trayectorias sintéticas, y hacer de manera continua pruebas de regresión con tareas en producción, para evitar que el modelo solo se desempeñe bien dentro de la distribución sintética.

\begin{warningbox}{Los tres grandes riesgos de los datos sintéticos}
\begin{itemize}
    \item Desvío semántico: la descripción del objetivo no coincide con la intención real de la tarea;
    \item Alucinación de acciones: se generan operaciones que parecen razonables pero no son ejecutables;
    \item Acumulación de sesgo: los datos generados por el modelo terminan reforzando sus propios patrones de error.
\end{itemize}
\end{warningbox}

\subsection{Resumen de la sección}
El uso correcto de los datos sintéticos es como «complemento y aceleración», no como sustituto de los datos reales de interacción. La experiencia clave que ofrece el curso es: convertir la verificación de la ejecución en una restricción dura del pipeline sintético.
\section{Diseño del modelo: arquitectura multietapa que unifica percepción, planificación y ejecución}

\subsection{Deficiencias clave de los GUI Agent actuales}
El expositor resumió tres tipos de problemas de los sistemas actuales: grounding visual inestable, planificación de largo alcance débil y mala generalización entre plataformas. Muchos modelos funcionan en tareas cortas, pero se degradan rápidamente en tareas de alta complejidad.

\begin{importantbox}{Tres tipos de carencias de capacidad}
\begin{enumerate}
    \item \textbf{Grounding}: percibe, pero no ubica con precisión, o hay ambigüedad al resolver el objetivo;
    \item \textbf{Planning}: ejecuta pasos individuales, pero carece de gestión global de subobjetivos;
    \item \textbf{Generalization}: el rendimiento cae abruptamente al cambiar de interfaz o de plantilla de tarea.
\end{enumerate}
\end{importantbox}

\subsection{El enfoque integrador que propone el curso}
Según los subtítulos, el expositor subraya que, mediante un paradigma generativo unificado de estilo AR, la comprensión visual, la cadena de razonamiento (monologue/reasoning trace) y la generación de acciones se incorporan a una misma interfaz de entrenamiento. Esto reduce las pérdidas de interfaz que produce el ensamblaje de múltiples modelos.

\begin{figure}[H]
\centering
\includegraphics[width=0.88\textwidth]{frame-07-model-architecture.jpg}
\caption{Video aproximadamente en 01:16:05: la arquitectura del modelo enfatiza un formato unificado de entrada y salida, así como objetivos de entrenamiento por etapas}
\end{figure}

\begin{knowledgebox}{Beneficios de ingeniería del entrenamiento por etapas}
\begin{itemize}
    \item Etapa 1: refuerza el grounding visual y la comprensión semántica de la interfaz;
    \item Etapa 2: incorpora la cadena de razonamiento y señales de planificación, para mejorar la toma de decisiones de largo alcance;
    \item Etapa 3: alinea la estrategia de ejecución con la retroalimentación de recompensa, para optimizar la tasa de finalización de tareas.
\end{itemize}
\end{knowledgebox}

\subsection{La incorporación de señales de planner}
En la parte final del curso se menciona repetidamente que el planner aporta beneficios significativos en tareas difíciles. Intuitivamente, el planner no sustituye al modelo de acción, sino que ofrece una estructura intermedia interpretable, que permite al modelo reducir la exploración a ciegas en tareas largas.

\begin{warningbox}{Planner no equivale a ``redundancia adicional''}
Si la tarea en sí tiene una cadena larga y un estado complejo, la ausencia de planner suele traducirse en más acciones inválidas y retrocesos. El costo adicional de la planificación suele ser menor que el costo de una ejecución errónea. La decisión de usar o no un planner debe tomarse por niveles, según la complejidad de la tarea.
\end{warningbox}

\subsection{Resumen de la sección}
Lo determinante en la arquitectura del modelo no es el tamaño de los parámetros, sino la combinación de ``interfaz unificada + entrenamiento por etapas + señales de planificación''. Estos tres elementos determinan en conjunto la estabilidad del sistema en tareas de alta complejidad.

\section{Resultados experimentales y conclusiones prácticas}

\subsection{Interpretación de los resultados: la tendencia importa más que los valores absolutos}
El curso muestra una tendencia de mejora en varios benchmarks, especialmente en tareas más difíciles y en configuraciones en línea, donde agregar planificación y una gobernanza más sólida de los datos de entrenamiento hace que el efecto sea más estable. Aunque las puntuaciones concretas varían según la configuración, la tendencia se mantiene muy consistente.

\begin{figure}[H]
\centering
\includegraphics[width=0.88\textwidth]{frame-08-results.jpg}
\caption{Video aproximadamente en 01:24:55: la página de resultados destaca cómo el entrenamiento por etapas y el planner elevan la tasa de éxito en tareas complejas}
\end{figure}

\begin{importantbox}{Tres lecciones de ingeniería reutilizables}
\begin{enumerate}
    \item Primero hay que estabilizar la evaluación y el entorno, y solo después hablar de la velocidad de iteración del modelo;
    \item la prioridad de la gobernanza de datos no es menor que la de las mejoras del modelo;
    \item para tareas complejas, la señal de planificación suele ser una ``condición necesaria'' y no ``un simple adorno''.
\end{enumerate}
\end{importantbox}

\subsection{Recomendaciones de métricas para la puesta en producción}
\begin{table}[H]
\centering
\begin{tabular}{p{3.2cm}p{5.0cm}p{2.8cm}}
\toprule
\textbf{Categoría de métrica} & \textbf{Métrica sugerida} & \textbf{Escenario de uso} \\
\midrule
Resultado de la tarea & Tasa de finalización, tasa de finalización parcial, distribución de tipos de fallo & Evaluar la utilidad del producto \\
Calidad de la trayectoria & Número promedio de pasos, tasa de acciones repetidas, tasa de retrocesos & Diagnosticar la calidad de la planificación \\
Estabilidad del sistema & Consistencia entre reejecuciones, tasa de recuperación tras interrupciones anómalas & Evaluar la confiabilidad en producción \\
Ciclo cerrado de datos & Proporción de muestras fallidas reincorporadas, magnitud de la mejora tras la corrección & Medir la eficiencia de la iteración \\
\bottomrule
\end{tabular}
\caption{Marco de métricas de ingeniería de agentes derivado de las conclusiones del curso}
\end{table}

\begin{warningbox}{No te fijes solo en la ``tasa de finalización total''}
Una sola métrica puede ocultar problemas. Por ejemplo, un aumento en la tasa de finalización puede deberse a que aumentó la proporción de tareas simples, y no a una mejora en la capacidad para tareas complejas. El curso recomienda hacer al menos un desglose estadístico por dificultad de la tarea, tipo de tarea y tipo de entorno.
\end{warningbox}

\subsection{Resumen de la sección}
Lo más valioso de los resultados de esta clase es la metodología: el entorno, la evaluación, los datos y el modelo deben avanzar de manera simultánea. Cualquier eslabón débil hará que el sistema se vuelva inestable en tareas reales.
\section{Hoja de ruta de implementación en ingeniería: del prototipo de investigación al sistema en producción}

\subsection{Etapa uno: plataforma de experimentación reproducible}
Primero se construye un entorno reproducible, un protocolo de acciones unificado y scripts de evaluación estandarizados. El objetivo de esta etapa no es perseguir la puntuación más alta, sino establecer una base de experimentación diagnosticable, regresionable y escalable.

\subsection{Etapa dos: ciclo de retroalimentación de datos}
Con las tareas fallidas reales como núcleo, se retroalimentan de manera continua a la construcción de datos y al proceso de entrenamiento. El impulso guiado por tutoriales y los datos sintéticos avanzan en paralelo, pero siempre manteniendo la restricción de verificación de ejecución, para evitar el desplazamiento de la distribución de los datos.

\subsection{Etapa tres: vigilancia en producción}
Se incorporan control de permisos, auditoría de operaciones, reversión ante anomalías y mecanismos de intervención manual. Solo cuando se cuenta con «observabilidad + capacidad de recuperación» puede el agente pasar de la demostración a la cadena crítica del negocio.

\begin{knowledgebox}{Requisitos de colaboración a nivel organizacional}
Un agente multimodal no es un proyecto que un solo equipo de modelos pueda completar por sí solo. Requiere la colaboración de ingeniería de entornos, ingeniería de evaluación, ingeniería de datos, entrenamiento de modelos y equipos de seguridad de producto, para formar un ritmo de lanzamiento unificado.
\end{knowledgebox}

\begin{importantbox}{Principio de implementación en una frase}
\textit{``Primero hay que ver con claridad el proceso de las fallas, y después buscar puntuaciones más altas.''} Esto aumenta la productividad real más que amontonar parámetros o muestras a ciegas.
\end{importantbox}

\subsection{Resumen de la sección}
La ruta de implementación debe seguir el principio de ``la plataforma primero, el ciclo de retroalimentación de datos, la vigilancia en producción''. El valor central del curso está en ofrecer una ruta de ingeniería de agentes ejecutable, escalable y auditable.

\section{Casos de fallo y manual de depuración: convertir un sistema inutilizable en un sistema iterable}

\subsection{Estratificación de tipos de fallo: primero agrupar, luego corregir}
En los proyectos de Agentes multimodales, el fallo no es la excepción sino la norma. Una idea de ingeniería clave implícita en el curso es: no fijarse directamente en el «fallo final», sino primero descomponer el fallo en subcategorías corregibles. Solo cuando los tipos de fallo se agrupan en categorías, las acciones de optimización resultan eficaces.

\begin{importantbox}{Clasificación recomendada de fallos en cuatro niveles}
\begin{enumerate}
    \item \textbf{Fallo de percepción}: no entender la interfaz, no leer el texto clave, ubicar incorrectamente el objeto objetivo;
    \item \textbf{Fallo de planeación}: orden incorrecto de subobjetivos, pasos omitidos, pasos repetidos en exceso;
    \item \textbf{Fallo de ejecución}: el formato de la acción es correcto pero no coincide con el contexto, o falla la llamada a la interfaz;
    \item \textbf{Fallo de verificación}: la tarea está casi completa, pero el script de verificación no coincide con la definición del objetivo.
\end{enumerate}
\end{importantbox}

\begin{knowledgebox}{Por qué «agrupar los fallos en categorías» afecta directamente la eficiencia del entrenamiento}
Si todas las muestras de fallo se reincorporan de manera uniforme, se amplifica el ruido. Una vez agrupadas por categorías, los datos pueden reforzarse de manera dirigida según el tipo de problema: para los fallos de percepción, muestras de grounding y OCR; para los fallos de planeación, muestras de razonamiento en cadenas largas; para los fallos de ejecución, muestras de restricción de acciones; para los fallos de verificación, muestras de prueba de robustez del script.
\end{knowledgebox}

\subsection{Ciclo cerrado de depuración: registros, reproducción y experimentos de control}
La práctica de ingeniería correspondiente al contenido del curso puede resumirse en tres pasos: primero, capturar el registro completo; segundo, reproducir la trayectoria; tercero, realizar un experimento de control mínimo. El objetivo del experimento de control no es obtener la mejor puntuación, sino confirmar si un cambio determinado realmente reduce cierto tipo de fallo.

\begin{table}[H]
\centering
\begin{tabular}{p{2.6cm}p{4.6cm}p{4.0cm}}
\toprule
\textbf{Etapa} & \textbf{Operación clave} & \textbf{Producto obtenido} \\
\midrule
Recolección de registros & Guardar observaciones, acciones, estado del entorno y salida del script de evaluación & Permite localizar el punto donde ocurre el fallo \\
Reproducción de trayectorias & Reproducir en orden temporal el comportamiento del Agente y los cambios de la interfaz & Atribución preliminar de la causa del fallo \\
Experimento de control & Sustitución de una sola variable: modelo, prompt, restricción de acciones o versión del script & Evidencia verificable de la corrección \\
\bottomrule
\end{tabular}
\caption{Flujo mínimo viable de depuración para Agentes multimodales}
\end{table}

\begin{warningbox}{Evitar «cambiar muchas cosas a la vez»}
Cuando el modelo, los datos, el formato de las acciones y el script de evaluación cambian al mismo tiempo, resulta casi imposible determinar el origen real de la mejora. El curso enfatiza la iteración ingenieril, que en esencia exige que cada ronda de experimentos sea explicable y reversible.
\end{warningbox}

\subsection{Ejemplos de problemas frecuentes}
\begin{itemize}
    \item \textbf{Desviación en la ubicación de elementos}: los cambios de escala de la interfaz provocan desviaciones de coordenadas; se recomienda introducir ubicación relativa y anclas de objeto;
    \item \textbf{Degradación en tareas largas}: la primera mitad es correcta y la segunda mitad falla; es común en la compresión de contexto y el olvido de la planeación;
    \item \textbf{Manipulación de la recompensa}: el Agente aprende a activar vulnerabilidades del script en lugar de completar la tarea, lo cual exige reforzar la robustez del script de verificación;
    \item \textbf{Desalineación de estado entre aplicaciones}: la ejecución paralela de varias ventanas y tareas en segundo plano provoca inconsistencias de estado, lo cual exige una gestión de sesión más sólida.
\end{itemize}

\begin{importantbox}{Priorización de la depuración}
Corregir primero los fallos «de alta frecuencia y alto costo», y después atender los problemas marginales de baja frecuencia. En un sistema de producción, la ganancia en estabilidad suele valer más que un desempeño óptimo en un solo caso.
\end{importantbox}

\subsection{Resumen de la sección}
La depuración de fallos no es un trabajo accesorio, sino el proceso principal de iteración del Agente. Solo al estructurar los fallos, hacerlos reproducibles y comparables, se forma un verdadero volante de entrenamiento.

\section{Agenda de investigación: tres frentes de avance para la siguiente generación de Agents multimodales}

\subsection{Frente de avance uno: capa de representación unificada}
El curso discutió extensamente las capturas de pantalla, los árboles de accesibilidad, el código de acciones y las cadenas de razonamiento, que en esencia son representaciones intermedias de distintas modalidades. Una dirección importante para la siguiente etapa es construir una capa de representación unificada que reduzca la pérdida de información al cambiar de modalidad.

\begin{knowledgebox}{Problemas que debe resolver la capa de representación unificada}
\begin{itemize}
    \item Cómo mapear los objetos visuales, los objetos semánticos y los objetos de acción a una misma estructura;
    \item Cómo lograr que distintas plataformas (web, escritorio, móvil) compartan una interfaz de abstracción;
    \item Cómo dar soporte a la compresión de estado y la recuperación reversible en tareas de cadena larga.
\end{itemize}
\end{knowledgebox}

\subsection{Frente de avance dos: planificación verificable de largo alcance}
Los modelos actuales todavía caen fácilmente en estrategias miopes ante tareas complejas. La parte final del curso enfatizó la contribución del planner, lo que indica que la «planificación verificable» será el próximo foco de competencia: no basta con generar un plan, sino verificar de manera continua si el plan sigue siendo coherente con el estado del entorno.

\begin{importantbox}{La evaluación de la capacidad de planificación debe pasar de «tiene o no tiene» a «calidad»}
Las evaluaciones futuras no solo deben preguntar si el modelo puede producir un plan, sino también evaluar la viabilidad de ejecución del plan, su capacidad de actualizarse y su eficiencia de recuperación ante errores. Sin estos indicadores, el planner puede reducirse a un texto intermedio meramente decorativo.
\end{importantbox}

\subsection{Frente de avance tres: seguridad y gobernanza nativas}
Cuando un Agent comienza a operar cuentas reales y sistemas de negocio, la seguridad ya no puede posponerse. Los límites de permisos, los registros de auditoría, la reversión ante anomalías y los mecanismos de intervención humana deben convertirse en restricciones nativas del entrenamiento y el despliegue del modelo, y no en parches previos al lanzamiento.

\begin{warningbox}{El mayor riesgo de un Agent no es solo «no poder hacerlo», sino «hacerlo mal y sin poder rastrearlo»}
Si un sistema carece de una cadena auditable, después de un incidente resulta imposible reconstruir los límites de responsabilidad. Aunque el curso se centra en la investigación, su pensamiento sobre entornos y evaluación ya ofrece una base práctica para la gobernanza nativa.
\end{warningbox}

\subsection{Lista de tareas ejecutables para los próximos 12 meses}
\begin{table}[H]
\centering
\begin{tabular}{p{2.8cm}p{4.8cm}p{3.6cm}}
\toprule
\textbf{Dirección} & \textbf{Tarea propuesta} & \textbf{Resultado esperado} \\
\midrule
Representación unificada & Construir una capa intermedia de objeto-acción entre plataformas & Reducir el costo de migración y la tasa de error en las acciones \\
Optimización de la planificación & Introducir verificación en línea del plan y mecanismos de replanificación & Aumentar la tasa de finalización de tareas largas \\
Volante de datos & Agrupar automáticamente las muestras de fallo por categorías y reincorporarlas al entrenamiento & Reducir la tasa de recurrencia de fallos frecuentes \\
Gobernanza de seguridad & Establecer plantillas de permisos, auditoría y estrategias de reversión & Sostener un despliegue en producción controlable \\
\bottomrule
\end{tabular}
\caption{Hoja de ruta que lleva la metodología del curso a la capa de ejecución del equipo}
\end{table}

\subsection{Resumen de la sección}
La competencia de la siguiente generación de Agents multimodales no dependerá solo de la escala de parámetros del modelo, sino de quién logre convertir la representación unificada, la planificación de largo alcance y los mecanismos de gobernanza en capacidades de ingeniería reutilizables.

\section{Síntesis y ampliación}

\subsection{Tabla de síntesis de toda la clase}
\begin{table}[H]
\centering
\begin{tabular}{p{3.0cm}p{4.6cm}p{4.0cm}}
\toprule
\textbf{Tema} & \textbf{Idea central} & \textbf{Acción de ingeniería} \\
\midrule
Definición del problema & El objetivo del agente pasa de responder preguntas al ciclo cerrado de tareas reales & Establecer una evaluación de interacción con el entorno, y no solo una calificación fuera de línea \\
Benchmark & OS-World llena el vacío de escenarios reales de computadora & Tareas guionizadas, entorno reproducible, resultados verificables \\
Sistema de datos & Los datos de trayectorias son el principal cuello de botella & Recolección impulsada por tutoriales + aumento sintético + gobernanza de calidad \\
Sistema del modelo & Unificar percepción, planeación y ejecución, y entrenar por etapas & Introducir señales de planeación, reforzar el grounding y la toma de decisiones de largo alcance \\
Estrategia de implementación & La viabilidad en producción depende de la observabilidad y de la capacidad de recuperación & Añadir mecanismos de auditoría, reversión, permisos y toma de control humano \\
\bottomrule
\end{tabular}
\caption{Abstracción de ingeniería de las conclusiones del curso}
\end{table}

\subsection{Lecturas adicionales}
\begin{itemize}
    \item El artículo y la documentación del proyecto OS-World (benchmark de agentes en entornos de computadora reales)
    \item WebArena / VisualWebArena / Mind2Web (la evolución de la evaluación de agentes web)
    \item Trabajos relacionados con el grounding de GUI y el modelado de acciones (de la localización visual a la ejecución de acciones)
    \item Artículos sobre la síntesis de datos de agentes multimodales y la selección de trayectorias (pipelines de datos de entrenamiento escalables)
    \item Prácticas de seguridad y auditoría de agentes (permisos, reversión, observabilidad)
\end{itemize}

\subsection{Preguntas de ampliación}
\begin{itemize}
    \item ¿Cómo diseñar una representación intermedia unificada para tareas entre aplicaciones que reduzca el acoplamiento entre plataformas?
    \item En un entorno de producción en línea, ¿cómo equilibrar la eficiencia de la automatización con los límites de seguridad?
    \item Cuando el objetivo de la tarea no puede guionizarse por completo, ¿cómo puede el sistema de evaluación coordinarse con la retroalimentación humana?
\end{itemize}

\end{document}
